\part{End-to-end validation on the 5G stack}\label{part:v32}

\begin{keybox}{What this part is}
The integrated closed loop was evaluated across all 32 class~$\times$~scenario combinations on Open5GS and UERANSIM,
under a predeclared plan (\code{docs/thesis/validation_32_plan.md}): 4 arms, 8 scenarios, 5 fresh seeds (50--54),
160 runs of 60\,s at 100\,Mb/s. The core signalling and GTP-U user plane are real; the radio is simulated. Results
are in \code{docs/VALIDATION_32.md}.
\end{keybox}

\section{Results}

\begin{qa}{V1}{How many of the 32 combinations pass?}
At the 0.95 SLA share: no slicing 19, rules 26, SliceNet with rules 24, the full system 23.
\end{qa}

\begin{qa}{V2}{Why do nine combinations fail under the full system?}
Four are uplink overload: healthcare and agriculture in disaster and medical emergency, where demand approaches
100\,Mb/s. Two are topology: education and general during a core switch failure, which has no alternate path. Three
appear only with SliceNet in the loop: general in fiber cut, and education and general in backhaul degradation. Those
three coincide with SliceNet's lowest flow accuracy, but the cause has not been diagnosed.
\end{qa}

\begin{qa}{V3}{What did the predeclared claims show?}
\begin{itemize}
\item Rules against no slicing, healthcare SLA: $+0.225$ $[0.130, 0.320]$, Holm $p<0.001$. Supported.
\item Full system against no slicing, worst class SLA: $+0.274$ $[0.192, 0.355]$, Holm $p<0.001$. Supported.
\item SliceNet against address rules, mean class SLA: $-0.007$ $[-0.015, 0.001]$, within the $-0.02$ margin.
  Supported.
\item RL against no RL, mean class SLA: $-0.001$ $[-0.006, 0.003]$. Unsupported.
\item RL against no RL, persistent ticks: $+2.7$ $[-0.4, 5.7]$. Unsupported.
\end{itemize}
Safety: 0 failed verifications and 0 validator rejections in 160 runs.
\end{qa}

\begin{qa}{V4}{Does that mean RL failed?}
RL ran correctly in the loop: all 17 runs with a persistent violation triggered it, and the safety gate accepted 81
of 173 proposals. It did not improve outcomes, and why has not been measured directly. One plausible reason is that
overload cannot be fixed by moving bandwidth within a full link. On this testbed, the evidence shows no gain from RL
over the rules allocator.
\end{qa}

\begin{qa}{V5}{Why was the first validation run discarded?}
A simulated radio link failed on the healthcare UE during seed 43. The gNB refused the reattach, and the tunnel kept
its address, so the readiness check passed and 56 more runs carried no healthcare traffic. The monitor also called
the dead slice healthy, because a dead path has no queue and no loss. Both defects were fixed. The plan requires
fresh seeds after a fix, so the matrix was repeated on seeds 50--54, and the void run is preserved.
\end{qa}

\begin{qa}{V6}{Is this a real 5G deployment?}
No. The 5G core and GTP-U user plane are real software, but the radio is UERANSIM, no radio resources are scheduled,
and it runs in Docker on one machine. Slicing acts on the transport path, not on the 5G core's policy functions.
\end{qa}
